% `tex.web` §1215 `get_r_token`, which every definition opens with:
%
%     begin restart: repeat get_token;
%     until cur_tok<>space_token;
%     if (cur_cs=0)or(cur_cs>frozen_control_sequence) then
%       begin print_err("Missing control sequence inserted");
%       help5("Please don't say `\def cs{...}', say `\def\cs{...}'.")
%       ("I've inserted an inaccessible control sequence so that your")
%       ("definition will be completed without mixing me up too badly.")
%       ("You can recover graciously from this error, if you're")
%       ("careful; see exercise 27.2 in The TeXbook.");
%       if cur_cs=0 then back_input;
%       cur_tok:=cs_token_flag+frozen_protection;
%       ins_error; goto restart;
%       end
%
% The recovery is the point.  The offending token goes *back*, a control
% sequence no document can name takes its place, and the definition runs to
% completion — so `\mathchardef A"7141` defines something harmless and then
% complains about `A` where a number was wanted, which is two errors and not
% one.  We returned nothing and every caller gave up in silence, leaving the
% rest of the definition to be read as ordinary text.
%
% §1224's `shorthand_def` opens the same way, so `\countdef`, `\chardef`,
% `\mathchardef`, `\toksdef` and the rest all reach it — that half was
% separately missing, since it asked only for a control sequence and took an
% active character for nothing.
\catcode`\{=1 \catcode`\}=2
\font\r=cmr10 \r \hsize100pt \parindent0pt
\mathchardef A"7141
\def B{x}
\countdef C=5
\let D=\relax
\chardef E=`\E
\toksdef F=3
\showthe\count5
\end
